\chapter*{Abstract}

This paper presents IncidentFlow, a domain-specific language (DSL) designed to improve cybersecurity incident reporting by reducing manual authoring effort and structural inconsistency. In current practice, incident reports are often produced through static templates and post-hoc editing workflows, where responders must repeatedly verify required sections, formatting constraints, and compliance-oriented evidence. As analyzed in Chapter 1, this process introduces procedural drift, weakens coordination visibility, and increases auditability risk when documentation is assembled under operational pressure.

The proposed solution models incident response and reporting within a single formal specification. Instead of reconstructing reports from fragmented notes, the DSL encodes triggers, phases, control flow, verification actions, severity transitions, and reporting events as first-class language constructs. This representation enables reporting requirements to be enforced by syntax and structure, thereby reducing repetitive requirement-checking by end users.

Implementation is based on an ANTLR-driven language pipeline, including a formally defined grammar, dedicated lexer rules, and parser logic with deterministic error handling, as detailed in Chapter 2. The resulting architecture provides a validated parse model that can support downstream automation of report generation and consistency checks.

The main contribution is a practical, extensible foundation for process-aware incident documentation that improves efficiency, reduces human error, and strengthens traceable reporting quality in high-pressure security operations.

\textbf{Keywords: } Domain-Specific Languages, Cybersecurity Incident Response, Incident Reporting Automation, ANTLR, Parser Design, Compliance Engineering, Human Error Reduction.
